\section{Multiscale Consistency (Restriction Inverse System and Inverse Limit)}\label{app:multiscale-inverse-limit}

This appendix completes the proofs from Section \ref{subsec:folding-multiscale}, making rigorous the mathematical meaning of ``the resolution $m$ does not alter the unified object; conclusions at different $m$ are consistent, differing only in the level of refinement.''

\begin{proof}[Proof of Lemma \ref{lem:pi-functorial-golden}]
For any $(x_1,\dots,x_{m_3})\in X_{m_3}$, both sides consist of truncating to the first $m_1$ digits, and hence are pointwise equal. The consistency of $\pi_{m_2}=\pi_{m_2}(\cdot)$ follows by the same reasoning: first truncating to the first $m_2$ digits and then to the first $m_1$ digits is the same as directly truncating to the first $m_1$ digits.
\end{proof}

\begin{proof}[Proof of Theorem \ref{thm:inverse-limit-golden}]
Define the map
\[
\Phi:X_\infty\to \varprojlim X_m,\qquad
\Phi(c):=\bigl(\pi_m(c)\bigr)_{m\ge 1}.
\]
By Lemma \ref{lem:pi-functorial-golden}, the family $(\pi_m(c))_{m\ge 1}$ satisfies the consistency condition of the inverse limit, so $\Phi$ is well-defined.

Conversely, given a consistent family $(x^{(m)})_{m\ge 1}\in\varprojlim X_m$, define the infinite sequence $c=(c_1,c_2,\dots)$ by setting $c_k:=x^{(m)}_k$ for any $m\ge k$. The consistency condition guarantees that this definition is independent of the choice of $m$. Moreover, each $x^{(m)}\in X_m$ satisfies the no-consecutive-$1$'s constraint, and hence $c\in X_\infty$.

Finally, we verify that these constructions are mutually inverse: for $c\in X_\infty$, the reverse construction returns $c$ by definition; for a consistent family $(x^{(m)})$, the reverse construction produces $c$ satisfying $\pi_m(c)=x^{(m)}$, thus returning the original consistent family. Therefore $\Phi$ is a bijection and the construction yields a natural isomorphism. \qed
\end{proof}
